\label{sec:prong2}

\subsection{The Legal Standard}

Prong 2 requires that the minority group be ``politically cohesive.''
\textit{Gingles}, 478 U.S. at 56.  Political cohesion means that minority
voters tend to support a common candidate in elections.  A minority group
that is divided in its political preferences does not satisfy Prong 2,
even if it is geographically concentrated.

Political cohesion has traditionally been demonstrated through racially
polarized voting (RPV) analysis: regression evidence that minority voters
support the same candidates at substantially higher rates than non-minority
voters.  The key quantity is the estimated minority vote share for the
minority-preferred candidate.

\subsection{WLS Regression for Minority Cohesion}

Minority political cohesion is estimated using the weighted least squares
(WLS) ecological regression introduced by \citet{goodman1953ecological} and
developed by \citet{king1997solution}.  The model is:

\begin{equation}
  Y_t = \alpha + \beta_{\text{min}} \cdot X^{\text{min}}_t
                + \beta_{\text{maj}} \cdot X^{\text{maj}}_t
                + \varepsilon_t
\end{equation}

where:
\begin{itemize}
  \item $Y_t$ is the precinct-level vote share for the minority-preferred
    candidate in precinct $t$.
  \item $X^{\text{min}}_t$ is the minority CVAP fraction in precinct $t$.
  \item $X^{\text{maj}}_t = 1 - X^{\text{min}}_t$ is the majority CVAP
    fraction.
  \item Weights are proportional to precinct total vote.
\end{itemize}

The estimator $\hat{\beta}_{\text{min}}$ is the estimated minority vote share
for the minority-preferred candidate---the Prong 2 quantity of interest.

\subsection{Primary Elections for Cohesion Analysis}

General election results conflate racial cohesion with partisan cohesion:
minority voters may appear cohesive in general elections simply because they
are predominantly aligned with one party.  \textit{Callais}~\citep{callais2026}
requires disentangling the two.  For Prong 2, the cleanest evidence of
racial cohesion independent of party comes from contested primary elections,
where candidates of the same party compete and party label does not determine
the result.

We use contested Democratic primary elections in districts with substantial
minority populations to estimate $\hat{\beta}_{\text{min}}$.  Where contested
primaries are unavailable, we use biracial general elections with a
partisan control variable (see Section~4 for the full disentanglement
methodology applicable to Prong 3).

\subsection{Running the Prong 2 Analysis}

The Prong 2 analysis is produced by:
\begin{verbatim}
redist analyze --state AL --year 2020 --version remedy \
  --types bloc-voting
\end{verbatim}
The \texttt{bloc-voting} analysis type runs the WLS regression and reports
$\hat{\beta}_{\text{min}}$ with HC3 standard errors and bootstrap confidence
intervals for each analyzed election.  The output identifies: (a) the
minority-preferred candidate in each election; (b) the estimated minority
vote share for that candidate; and (c) the statistical significance of the
minority cohesion estimate.

A standard Prong 2 evidentiary package includes at least three elections
spanning at least two election cycles, demonstrating that minority cohesion
is a persistent pattern rather than a single-election artifact.
